\documentclass[a4paper,11pt]{article}
\usepackage[utf8]{inputenc}
\usepackage[T1]{fontenc}
\usepackage[french]{babel}
\usepackage{geometry}
\geometry{left=2cm,right=2cm,top=1.5cm,bottom=1.5cm}
\usepackage{hyperref}
\usepackage{enumitem}

\begin{document}

\begin{flushleft}
    \textbf{AISSAOUI Ismail} \\
    Ingénieur en Intelligence Artificielle \& Data Science \\
    Chlef, Algérie \\
    Tél : +213 660 70 77 96 \\
    Email : \href{mailto:ismail.aissaoui.pro@gmail.com}{ismail.aissaoui.pro@gmail.com} \\
    Portfolio : \href{https://ismail-aissaoui.vercel.app}{ismail-aissaoui.vercel.app}
\end{flushleft}

\begin{flushright}
    À l'attention de \\
    \textbf{M. David Panzoli} (\href{mailto:david.panzoli@univ-jfc.fr}{david.panzoli@univ-jfc.fr}) \\
    \textbf{M. Ludovic Hamon} (\href{mailto:ludovic.hamon@univ-lemans.fr}{ludovic.hamon@univ-lemans.fr}) \\
    PEPR eNSEMBLE \\
    Fait le \today
\end{flushright}

\vspace{0.5cm}

\noindent \textbf{Objet : Candidature au contrat doctoral : « Modèle de co-apprentissage médié par un tuteur IA dans les formations immersives »}

\vspace{0.5cm}

Messieurs,

\medskip

C'est avec un vif intérêt que je vous adresse ma candidature pour le sujet de thèse intitulé « Modèle de co-apprentissage médié par un tuteur IA dans les formations immersives », au sein du projet PEPR eNSEMBLE. Étudiant en dernière année d'ingéniorat en Intelligence Artificielle à l'École Supérieure en Informatique (ESI-SBA), je suis passionné par l'intersection entre l'IA adaptative et les environnements de simulation complexes.

Mon parcours académique et mes expériences de recherche m'ont permis de développer une expertise robuste dans les domaines clés de votre appel à candidature. Dans le cadre de mon stage chez \textbf{Sonatrach}, j'ai conçu un \textbf{Jumeau Numérique} pour des turbines à gaz, impliquant la modélisation de flux physiques et l'estimation de la durée de vie résiduelle (RUL) par Deep Learning. Cette immersion dans la simulation de haute fidélité m'a sensibilisé aux enjeux de la réplication numérique de processus critiques, un socle technologique indispensable pour les environnements de formation immersifs.

Ma maîtrise des architectures \textbf{Transformers} et du \textbf{Reinforcement Learning} (RL), illustrée par le développement du système « Aura » (DQN/LSTM), constitue un atout majeur pour concevoir le tuteur IA incarné. Je suis particulièrement stimulé par la problématique du co-apprentissage écologique : je maîtrise des concepts tels que l'\textbf{Imitation Learning} et l'\textbf{Inverse Reinforcement Learning} (IRL), qui sont des pistes prometteuses pour permettre à un tuteur d'apprendre du geste de l'expert tout en s'adaptant à la progression de l'apprenant. Mes travaux récents sur la reconnaissance d'actions et de conformité (YOLOv8) renforcent ma capacité à modéliser et évaluer le \textbf{geste technique} avec précision.

Bien que mes projets récents utilisent principalement Python et les frameworks web/desktop, ma formation d'ingénieur m'a doté d'une excellente rigueur en C++ et d'une grande agilité pour maîtriser des moteurs comme \textbf{Unity3D} ou \textbf{Unreal Engine}. Ma capacité à traduire des concepts théoriques (Symbolic AI, Graph Learning) en démonstrateurs techniques performants sera un levier pour la réussite de ce projet doctoral.

Le prestige du PEPR eNSEMBLE et la nature interdisciplinaire de cette thèse (XR, IA, Didactique) correspondent parfaitement à mon ambition de chercheur. Je suis convaincu que mon profil d'ingénieur spécialisé en IA apportera une valeur ajoutée significative à la production de contributions méthodologiques et techniques validant votre approche de co-apprentissage.

Vous trouverez ci-joint mon CV détaillé ainsi que mes relevés de notes. Je reste à votre entière disposition pour un entretien afin de discuter plus en détail de mon adéquation avec votre projet doctoral.

Je vous prie d'agréer, Messieurs, l'expression de mes salutations distinguées.

\vspace{0.4cm}

\begin{flushleft}
    \textbf{AISSAOUI Ismail}
\end{flushleft}

\end{document}
